\chapter{Discussion}
\label{chap:discussion}

\section{Interpretation of the Combined Evidence}

This study began with a plausible forecasting proposition: a compact latent macroeconomic state might improve loan-level default probabilities by changing both the common level of risk and selected borrower sensitivities. The terminal split initially supported that proposition. The regime-aware model had the lowest Brier score and log loss and the highest ROC-AUC and PR-AUC in April 2023--September 2024. Those results remain correct for that fitted pipeline and period. They are not, however, sufficient evidence that regime awareness produces a recurring forecasting advantage.

The causal rolling-origin analysis changes the weight assigned to the terminal result. Each origin estimated macroeconomic transformations, scaling parameters, the HMM, logistic coefficients, regularisation strength, and Platt calibration without test-period outcomes. It then repeated the comparison over three successive and non-overlapping twelve-month tests. Regime-aware was best in the first test and worse than borrower-only in the next two. When the 36 test months were pooled, borrower-only was better on both proper scores, both discrimination measures, and monthly default-count calibration. The point estimates therefore support borrower-only as the more stable forecasting specification in this dataset.

The paired moving-block analysis adds an important qualification. Its intervals included zero for both the Brier and log-loss differences. This does not overturn the pooled ordering, and it does not establish equivalence. It indicates that the effective temporal evidence is much smaller than the count of loan-month rows might suggest. Millions of loans share only 36 rolling-test calendar months. Resampling short monthly blocks accordingly produced sequences in which the period-specific regime-aware advantage received different weight. The defensible result is a pooled advantage for borrower-only accompanied by uncertainty about the exact magnitude of the difference.

The combined evidence is thus more informative than either evaluation alone. The terminal split demonstrates that the regime-aware structure can be useful in a particular period. The rolling analysis demonstrates that this usefulness did not persist automatically when the forecast origin moved. Together they support a period-dependent diagnostic interpretation rather than a universal model-win claim.

\section{A Calibration-First Reading}

Probability forecasts should first be assessed as probabilities. In the rolling tests, borrower-only had a pooled Brier score of 0.005394 and log loss of 0.031469, compared with 0.005481 and 0.033326 for regime-aware. These proper scores favour borrower-only. The ordering is reinforced by monthly count calibration: borrower-only had a mean absolute default-count error of 1,009.7, while regime-aware had 1,266.3. Regime-aware had lower monthly loss in only 15 of 36 months.

Neither principal model was ideally calibrated in the large. The pooled event rate was 0.005538, while mean predicted PD was 0.011004 for borrower-only and 0.012392 for regime-aware. The validation-fitted Platt maps therefore did not fully transport to the following test periods. This common overprediction matters operationally. A model can rank individual loans effectively while overstating expected portfolio defaults, capital needs, or the volume of accounts entering an intervention queue. The monthly aggregation makes that consequence visible.

The terminal split illustrated a related distinction. Regime-aware improved pooled Brier score and ranking there, but borrower-only produced the smaller average monthly count error. The apparent inconsistency is only an aggregation difference. A probability model can reallocate risk more effectively among loans and still predict too many events in aggregate. For this reason, a small improvement in a pooled loan-level score should not be presented as a complete calibration success.

ROC-AUC and PR-AUC remain useful because defaults are rare and ordering matters for review or intervention. Yet they do not repair probability levels. In the pooled rolling tests, borrower-only was also stronger on these measures, with ROC-AUC 0.784796 and PR-AUC 0.103932 compared with 0.753599 and 0.089320 for regime-aware. The agreement of proper-score, ranking, and count-calibration evidence makes the primary rolling conclusion clearer than it would be if the metrics disagreed.

The additive macro model provides a warning about temporal calibration. Its first-test mean PD of 18.6756\% followed a validation period with a much higher event rate than the subsequent test. A validation-fitted monotone calibration map could not guarantee transport when the economic and portfolio distribution changed. This failure concerns the retained additive specification and its calibration under the observed shift. It is not evidence that macroeconomic information is universally unhelpful.

\section{Meaning and Limits of the Regime Signal}

The two-state HMM fulfilled its intended compression role. It mapped four adverse-oriented macroeconomic factors into a filtered probability, and its transition structure represented persistence rather than classifying each month independently. The regime-aware logit then supplied an explicit pathway through which the common state could shift loan-level log odds and alter eight selected sensitivities. This is substantively more interpretable than an unrestricted interaction search.

The estimated posterior nevertheless approached a hard classifier. In the terminal test every month had stress probability above 0.999902. Under causal rolling fitting, 35 of 36 test months had probability at or beyond 0.001 and 0.999. The second rolling test was entirely stress, and the first contained only one hard-stress month. Those values leave little support for interpreting a smooth change from normal to stress inside a fold.

Saturation does not make the model scores invalid. Predictions can still differ across loans because borrower characteristics vary and because interaction coefficients were estimated before the test. It does restrict the economic claim. The analysis cannot reliably estimate how forecast performance changes across intermediate stress probabilities when almost no intermediate probabilities occur. Hard-state comparisons within an all-stress or almost-all-normal test are similarly weak. The HMM is best read here as an origin-specific period classification and an interpretable model input, not as a precisely measured continuous stress intensity.

The independently fitted rolling HMMs also show parameter sensitivity. State labels were consistently oriented by the adverse-factor rule, but their test-state composition changed with the training window. This is expected in a latent-variable model: numeric labels have no inherent meaning, and emission and transition estimates depend on the history available at the origin. Deterministic multiple starts reduced optimisation dependence, while training likelihood selected among valid starts without test information. It did not remove sampling uncertainty from a model estimated on a relatively short monthly series.

The saturated posterior is therefore both a result and a limitation. It indicates that the fitted Gaussian states were sharply separated under the selected factors. It also means that a hard threshold of 0.5 contributed little additional information over most evaluation windows. A richer state model, alternative covariance restriction, or different factor construction might produce a smoother probability, but choosing such variants after observing these test results would constitute new model search. They are appropriate subjects for future research rather than retrospective changes to the frozen thesis comparison.

\section{Interpretability and Predictive Stability}

Interpretability is a property of model structure, not proof of forecast quality. The regime-aware equation identifies the borrower variables whose log-odds sensitivities are permitted to change with $q_t$. It also separates that channel from a common stress-level term. This makes the source of a probability change inspectable and provides a coherent diagnostic language for discussing borrower vulnerability under fitted adverse conditions.

The rolling results show why transparency and stability must be reported separately. A model can have a clear equation and still transport poorly when the macroeconomic state, portfolio composition, event prevalence, or calibration relationship changes. Conversely, the simpler borrower-only model can omit an explicit macroeconomic channel and still be the stronger pooled forecast because behavioural and contract variables already capture substantial cross-sectional risk.

The regime-aware model should consequently not replace borrower-only solely because it has a richer economic interpretation. A defensible use is parallel monitoring: borrower-only supplies the primary probability forecast, while regime-aware supplies an interpretable sensitivity diagnostic and a challenger prediction. Material divergence between them can prompt review of macroeconomic conditions, calibration drift, and borrower segments. That use respects the observed period dependence rather than hiding it behind one pooled terminal score.

\section{Relationship to Existing Work}

The results are consistent with the broad literature in a qualified way. Hidden-state and regime-switching credit models establish that persistent economic conditions can be represented explicitly in migration, intensity, and structural-risk settings \citep{bangia_et_al_2002,yu2018modeling,milidonis_chisholm_2024}. Loan-level macroeconomic studies likewise provide reasons to expect unemployment, interest rates, house prices, and related conditions to affect credit outcomes \citep{bellotti2009credit,bellotti2013stress}. The present study confirms that an interpretable regime channel can add predictive value in some periods.

It also supplies a stricter caution. A terminal holdout can favour a regime-aware model even when successive causal origins do not. Full-window scaling and HMM estimation answer a narrower retrospective parameter-estimation question than origin-specific fitting. By moving every upstream transformation inside the training window and retaining a fixed model family, the rolling pass tests whether the apparent gain survives a forecast-like sequence. In this study it did not recur.

This conclusion does not contradict the conceptual value of regime switching. It distinguishes a mechanism from a performance guarantee. A latent state can be economically meaningful and diagnostically useful without producing lower forecast loss in every period. The result also accords with the forecasting principle that model comparisons depend on horizon, sample, information set, calibration method, and evaluation dates. Claims should remain conditional on those choices.

No unreviewed preprint is used as a core authority in this thesis. In particular, a recent preprint describing a different gradient-boosting and regime-specific calibration system was excluded from the bibliography pending independent verification and supervisor guidance. Its model class and uniformly favourable reported outcome do not provide a like-for-like benchmark for the fixed logistic comparison conducted here.

\section{Practical Implications}

For portfolio forecasting, the evidence favours borrower-only logistic regression as the primary model among the three frozen candidates. It delivered the best pooled causal proper scores, discrimination, and monthly count calibration. This recommendation is comparative, not absolute: its mean PD still exceeded the rolling event rate, so deployment would require continuing calibration surveillance.

The regime-aware model remains useful in three narrower roles. First, it can act as a challenger that indicates whether a macro-conditioned specification is beginning to outperform the primary model. Second, its interaction structure can support qualitative review of which loan characteristics are allowed to become more or less influential under the fitted state. Third, the HMM path can provide a compact descriptive indicator for scenario discussion. None of these roles requires claiming that the model is universally more accurate.

Operational use should preserve the information boundaries of the rolling study. Scaling, HMM parameters, logistic coefficients, and calibration maps should be refitted on a documented schedule using only information available at the forecast date. Monitoring should include Brier score, log loss, mean predicted PD versus event rate, grouped calibration, monthly expected versus realised default counts, and both ROC-AUC and PR-AUC. Large changes in the selected HMM start, state composition, or posterior entropy should be treated as diagnostic events rather than silently absorbed.

Threshold-based credit decisions were outside the scope of this study. The reported metrics evaluate probability forecasts and rankings, not approval, pricing, collection, or capital thresholds. Any decision threshold would require an explicit loss function, capacity constraint, fairness and governance review, and validation on the intended decision population.

\section{External Validity and African Context}

The empirical data describe United States mortgages, not African consumer or mortgage portfolios. Freddie Mac loans are affected by United States underwriting conventions, servicing rules, interest-rate transmission, labour markets, housing markets, and public data infrastructure. The FRED variables and their transformations likewise encode a United States macroeconomic environment. Coefficients, event rates, regime dates, and calibration maps cannot be transferred directly to an African lender.

The relevance to an African research audience is therefore methodological rather than geographic. The study offers a transparent pipeline for combining loan-level behavioural data with a causal filtered macroeconomic state, and it demonstrates why calibration-first rolling evaluation can overturn a favourable terminal result. This lesson is pertinent wherever macroeconomic variation is substantial and usable credit histories are temporally structured. It should be presented as a validation framework and cautionary empirical case, not as evidence about African borrower behaviour.

\section{Limitations}

First, the loan data are sample files rather than the full Freddie Mac population, and the study considers one twelve-month 90-DPD transition definition. Results may differ for another product, event threshold, horizon, cohort range, or risk-set rule. Second, current-vintage FRED observations were aligned historically. The design prevents future dates from entering rolling transformations but does not reconstruct the real-time vintage that an analyst would have observed at every origin.

Third, the macroeconomic effective sample is monthly and short relative to the loan-month panel. HMM and block-uncertainty conclusions are consequently supported by far fewer independent temporal units than the raw observation count suggests. Fourth, the Gaussian two-state full-covariance HMM is deliberately fixed and parsimonious. Posterior saturation limits smooth-regime interpretation, while alternative state counts or emission families were not searched after test inspection.

Fifth, calibration was estimated on the immediately preceding validation block and applied unchanged to each test. This enforces honest separation but exposes the forecasts to prevalence and relationship shifts. Platt scaling was retained as the primary method by protocol; isotonic results from the terminal sensitivity analysis were not used to choose the rolling method. Sixth, the paired moving-block intervals operate on monthly average losses. They preserve short temporal clustering but do not jointly resample loans and months or quantify every source of parameter uncertainty.

Finally, the study is predictive rather than causal. Coefficients and interactions are conditional associations in a regularised model. The HMM state is a fitted statistical summary, not an observed economic treatment. These limits do not remove the comparative evidence; they determine how narrowly it should be stated.
